%! AP Statistics — Unit 4, Lesson 4.1 — Warm-Up
\documentclass[10pt]{article}
\usepackage{apstats-article}
\usepackage{apstats-boxes}

\begin{document}

\pageheader{Unit 4, Lesson 4.1}{Warm-Up}
\namedateperiod

\vspace{0.3in}

% ── SECTION 1: Spiral Review ──────────────────────────────────────────────────
{\large\bfseries\color{navy} Part 1 \quad Spiral Review}

\vspace{0.12in}

\begin{enumerate}

    \item From Unit 3 --- classify each data collection method as \textbf{random} (R)
    or \textbf{non-random} (NR), then state whether it produces trustworthy conclusions.

    \vspace{8pt}
    \renewcommand{\arraystretch}{1.8}
    \begin{tabularx}{\linewidth}{@{} X p{2cm} p{3.5cm} @{}}
    \textbf{Method} & \textbf{R or NR?} & \textbf{Trustworthy?} \\
    \hline
    A researcher randomly dials 500 phone numbers. &
    \blank{1.5cm} & \blank{3cm} \\
    A newspaper invites readers to email their opinion. &
    \blank{1.5cm} & \blank{3cm} \\
    A teacher surveys her own students to estimate school-wide opinion. &
    \blank{1.5cm} & \blank{3cm} \\
    \end{tabularx}

    \vspace{0.18in}

    \item A scatterplot shows that students who spend more hours per week on social
    media tend to score lower on reading tests ($r = -0.62$).
    \begin{enumerate}[label=(\alph*), itemsep=4pt]
        \item Describe the direction, form, and strength of the association.\\
        \writeline
        \item Can we conclude that social media use \emph{causes} lower reading scores?
              Explain.\\
        \writeline\writeline
    \end{enumerate}

\end{enumerate}

\vspace{0.35in}

% ── SECTION 2: Looking Ahead ──────────────────────────────────────────────────
{\large\bfseries\color{navy} Part 2 \quad Looking Ahead}

\vspace{0.12in}

\noindent Unit 4 asks: \textit{``Could this have happened by chance?''}

\vspace{0.14in}

\begin{tcolorbox}[colback=greenbg, colframe=greenacc, boxrule=0.8pt, arc=2mm,
    left=4mm, right=4mm, top=3mm, bottom=3mm]
    \small
    \begin{itemize}[leftmargin=1.3em, itemsep=8pt]
  \item A random process is any process whose outcome is determined by chance.
  \item Patterns in data do \textbf{not} necessarily mean the variation is non-random.
  \item Probability is the tool that lets us decide whether a pattern is surprising or expected.
    \end{itemize}
\end{tcolorbox}

\vspace{0.2in}

\noindent\textcolor{slate}{\small Think about a time you saw a streak or pattern (in sports, weather,
a game, etc.) and assumed it meant something. Write your example below:}\\[8pt]
\writeline\\[2pt]\writeline

\end{document}
